Локализация и позиционирование -- ключевые элементы системы навигации и управления РНТТМ, позволяющий определить их положение и~движение в пространстве. Для определения положения, ориентации, скорости и ускорения РНТТМ могут использоваться различные данные. Как уже было отмечено на рисунке~\cref{fig:arch_tech}, основные компоненты включают:

\begin{itemize}
	\item Global Navigation Satellite System (GNSS)	-- система спутниковой навигации, которая определяет местоположение на поверхности Земли. Она объединяет данные от нескольких спутниковых систем, таких как GPS (США), ГЛОНАСС (РФ), Galileo (EC) и BeiDou (КНР). Предоставляет данные о местоположении, но может быть неточным в условиях плохой видимости спутников, например, в городах или тоннелях;

	\nomenclature{GNSS}{Global Navigation Satellite System\nomrefpage}
	\nomenclature{США}{Соединённые Штаты Америки\nomrefpage}
	\nomenclature{РФ}{Российская Федерация\nomrefpage}
	\nomenclature{ЕС}{Европейский Союз\nomrefpage}
	\nomenclature{КНР}{Китайская Народная Республика\nomrefpage}

 	\item инерциальный измерительный модуль -- устройство, состоящее из акселерометра и гироскопа. Оно измеряет ускорение и угловую скорость, что помогает отслеживать движение робота. Его показания могут искажаться при вибрациях или наклонах;

	\item датчики вращения колес — высокоточные устройства, установленные на колесах. Они фиксируют повороты колес, что позволяет точно измерять пройденное расстояние и скорость. Точность снижается на скользких или неровных поверхностях;

	\item высококачественная картография -- карты высокой детализации, содержащие трехмерные данные об окружающей среде. Они состоят из лидарного и векторного слоев:

	\begin{enumerate}[left=0pt]
		\item \textit{лидарный слой}:
		трехмерная модель, созданная на основе данных лидара. Она формируется автоматически: РНТТМ проезжает по местности, собирает данные, которые затем обрабатываются для создания детальной карты. Динамические объекты, такие как машины или пешеходы, должны удалятся из данных, чтобы карта оставалась стабильной. Для этого может применяется алгоритм Normal Distribution Transform, который сравнивает текущие данные лидара с картой;

		\item \textit{векторный слой}:
		часть высококачественной картографии, содержащая информацию о дорожной инфраструктуре. Этот слой включает данные о пешеходных переходах, тротуарах, велодорожках и других зонах. Также в нем хранится информация о светофорах и их взаимосвязи, что помогает роботу планировать маршрут.
	\end{enumerate}
\end{itemize}

В направлении объединения этих данных имеются известные наработки, например метод Extended Kalman, который соединяет данные от всех сенсоров для получения точного результата. Он учитывает погрешности каждого сенсора и корректирует данные, чтобы минимизировать ошибки. Однако для его работы требуется правильная настройка, иначе могут возникать неточности.

Таким образом, локализация РНТТМ достигается за счет комбинации данных от различных сенсоров и алгоритмов, что позволяет точно определять его положение и движение в пространстве.